% Fig — detailed process flow of the flux-anchored inner loop: Step A (settle the
% skin) and Step B (reconstruct the deep). All values from src/lunar (config.py,
% equilibrium.py, solver.py). Body only; standalone wrapper supplies preamble.
\begin{tikzpicture}[
  font=\footnotesize,
  every node/.style={align=center},
  box/.style={rectangle, rounded corners=2pt, draw=dim, line width=0.8pt,
              fill=white, inner sep=5pt, text width=7.1cm},
  io/.style={rectangle, rounded corners=10pt, draw=dim, line width=0.9pt,
             fill=black!3, inner sep=6pt, text width=7.1cm},
  aBox/.style={box, draw=coral},
  bBox/.style={box, draw=forest},
  test/.style={rectangle, rounded corners=2pt, draw=teal, line width=1pt,
               fill=teal!8, inner sep=5pt, text width=6.0cm},
  arrow/.style={-{Stealth[length=2.4mm]}, line width=0.9pt, draw=dim,
                rounded corners=3pt},
  steplbl/.style={anchor=south west, font=\small\bfseries}]

  % ---- spine nodes ----
  \node[io] (start)
    {\textbf{Start.}  Flat guess $T_\mathrm{guess}$; truncate the grid to
     $[0,\;z_0+0.15\,\mathrm{m}]$ so only the skin is time-stepped.};
  \node[box, below=5mm of start] (sched)
    {\textbf{Two-stage anchor schedule}\\[2pt]
     stage 1:\ \ $z_0=0.25$ m,\ \ $n_\mathrm{in}=4$,\ \ tol $=0.10$ K\\
     stage 2:\ \ $z_0=0.55$ m,\ \ $n_\mathrm{in}=96$,\ \ tol $=0.005$ K};

  \node[aBox, below=11mm of sched] (a1)
    {March $n_\mathrm{in}$ lunations of Crank--Nicolson ($\Delta t=1800$ s
     $+$ wrap step, 2 Picard sweeps,
     $\theta=\tfrac12$) on the truncated grid --- nonlinear surface energy balance
     (Newton), harmonic-mean face $K$, Thomas tridiagonal solve --- until the
     diurnal skin is periodic.};
  \node[aBox, below=5mm of a1] (a2)
    {$\Rightarrow$ cycle-mean skin $\langle T\rangle(z)$ for $z\le z_0$, the anchor
     value $\langle T\rangle(z_0)$, and the saved sub-daily flux history.};

  \node[bBox, below=12mm of a2] (b1)
    {Rectified (eddy) flux from that history:\\[3pt]
     $u_\mathrm{rect}=\langle K\,\partial_z T\rangle-K(\langle T\rangle)\,
      \partial_z\langle T\rangle$\\[2pt]
     {\itshape\color{dim}$<1\%$ of $Q_b$ below the anchor.}};
  \node[bBox, below=5mm of b1] (b2)
    {RK2 (midpoint) integrate the closure ODE {\itshape downward} from $z_0$:\\[4pt]
     $\dfrac{d\langle T\rangle}{dz}=\dfrac{Q_b-u_\mathrm{rect}}{K(\langle T\rangle,\,z)}$
     \\[4pt]
     $\Rightarrow$ full column $\langle T\rangle(z)$ to 5 m --- the deep part is
     reconstructed, never time-stepped.};

  \node[test, below=10mm of b2] (dec)
    {anchor drift\\ $|\Delta\langle T\rangle(z_0)| < \text{tol}$ ?};
  \node[io, below=10mm of dec] (done)
    {\textbf{Certified steady} $\langle T\rangle(z)$: passes guess-independence
     ($\le0.05$ K), flux closure, and the 120-lunation drift re-run ($\le0.08$ K).
     Feeds the $K_d$ RMSE.};

  % ---- spine arrows ----
  \draw[arrow] (start) -- (sched);
  \draw[arrow] (sched) -- (a1);
  \draw[arrow] (a1) -- (a2);
  \draw[arrow] (a2) -- (b1);
  \draw[arrow] (b1) -- (b2);
  \draw[arrow] (b2) -- (dec);
  \draw[arrow] (dec) -- (done)
    node[midway, right=1pt, font=\scriptsize\color{dim}] {yes};

  % ---- loop-back (no): tight left channel back to the top of Step A ----
  \path (dec.west) -- ++(-1.25,0) coordinate (turn);
  \draw[arrow] (dec.west) -- node[above=1pt, font=\scriptsize\color{dim}] {no}
               (turn) |- (a1.west);

  % ---- step-group containers + labels ----
  \begin{scope}[on background layer]
    \node[fill=coral!7,  rounded corners=3pt, fit=(a1)(a2), inner sep=3mm] (contA) {};
    \node[fill=forest!7, rounded corners=3pt, fit=(b1)(b2), inner sep=3mm] (contB) {};
  \end{scope}
  \node[steplbl, color=coral]  at ([yshift=1.5mm]contA.north west)
    {STEP A \textbullet\ settle the skin};
  \node[steplbl, color=forest] at ([yshift=1.5mm]contB.north west)
    {STEP B \textbullet\ reconstruct the deep};

\end{tikzpicture}
